% chktex-file 36
\section{Chức năng cp (Sao chép tập tin)}

\subsection{Mục tiêu}
Chức năng \texttt{cp} cung cấp tiện ích sao chép tệp tin cơ bản trên hệ thống tệp của xv6 theo cú pháp \texttt{cp src dst}. Chương trình đọc nội dung từ tệp nguồn \texttt{src} và ghi sang tệp đích \texttt{dst} (tạo mới nếu chưa tồn tại hoặc ghi đè nếu đã có sẵn). Chương trình có cơ chế xử lý lỗi chặt chẽ: thông báo lỗi khi thiếu đối số, không thể mở tệp nguồn hoặc không thể tạo/ghi tệp đích.

\subsection{Tổ chức triển khai}
Chương trình được cài đặt trong tệp \texttt{user/cp.c}. Danh sách \texttt{UPROGS} trong \texttt{Makefile} được bổ sung mục \texttt{user/\_cp} để đưa tệp thực thi vào hệ thống tệp xv6.

Chương trình sử dụng các system call và hằng số cờ của hệ điều hành:
\begin{itemize}
  \item \texttt{open(path, flags)}: Mở tệp nguồn với cờ \texttt{O\_RDONLY} (chỉ đọc). Mở hoặc tạo tệp đích với tổ hợp cờ \texttt{O\_WRONLY | O\_CREATE | O\_TRUNC} (mở để ghi, tạo mới nếu chưa có, hoặc xóa trắng nội dung cũ để ghi đè).
  \item \texttt{read()} và \texttt{write()}: Thực hiện truyền dữ liệu giữa hai tệp thông qua vùng đệm \texttt{char buf[512]}.
  \item \texttt{fprintf(2, ...)}: Xuất các thông báo lỗi trực tiếp vào luồng xuất lỗi chuẩn (\texttt{stderr}, file descriptor 2).
  \item \texttt{close(fd)}: Đóng các tệp mô tả sau khi hoàn tất hoặc khi phát sinh sự cố.
\end{itemize}

\subsection{Luồng xử lý}
\begin{enumerate}
  \item \textbf{Kiểm tra tham số:} Kiểm tra số lượng đối số dòng lệnh \texttt{argc}. Nếu $\texttt{argc} \neq 3$, chương trình xuất thông báo \texttt{usage: cp src dst} ra \texttt{stderr} và thoát với mã lỗi 1 (\texttt{exit(1)}).
  \item \textbf{Mở tệp nguồn:} Gọi \texttt{open(argv[1], O\_RDONLY)}. Nếu trả về giá trị âm (không thể mở tệp, tệp không tồn tại), in \texttt{cp: cannot open <src>} và thoát với mã 1.
  \item \textbf{Mở tệp đích:} Gọi \texttt{open(argv[2], O\_WRONLY | O\_CREATE | O\_TRUNC)}. Nếu thất bại, in thông báo lỗi, đóng tệp nguồn đã mở trước đó và thoát với mã 1.
  \item \textbf{Sao chép dữ liệu:} Sử dụng vòng lặp \texttt{while} với hàm \texttt{read(fd\_src, buf, sizeof(buf))}:
    \begin{itemize}
      \item Mỗi lần lặp đọc tối đa 512 bytes vào bộ đệm \texttt{buf}.
      \item Ghi đúng số byte vừa đọc được sang tệp đích bằng hàm \texttt{write(fd\_dst, buf, n)}. Nếu số byte ghi được khác \texttt{n}, báo lỗi \texttt{cp: write error}, đóng cả hai tệp và thoát.
    \end{itemize}
  \item \textbf{Kết thúc:} Khi \texttt{read()} trả về 0 (đã đọc hết tệp - EOF), chương trình đóng cả hai tệp mô tả bằng \texttt{close()} và kết thúc thành công bằng \texttt{exit(0)}.
\end{enumerate}

\subsection{Khó khăn và giới hạn}
\begin{itemize}
  \item \textbf{Tối ưu hóa I/O bằng Buffer:} Thay vì sao chép từng byte một (gây ra quá nhiều lần chuyển ngữ cảnh \textit{context switch} giữa User mode và Kernel mode), chương trình sử dụng khối đệm 512 bytes. Kích thước này tương thích tốt với kích thước sector/block đĩa của xv6, giúp tối ưu hóa số lần gọi system call và tăng tốc độ sao chép.
  \item \textbf{Đảm bảo giải phóng tài nguyên khi có lỗi:} Khi quá trình mở tệp đích hoặc ghi dữ liệu thất bại giữa chừng, chương trình luôn đóng các file descriptor đã mở trước khi gọi \texttt{exit(1)} để tránh rò rỉ tài nguyên hệ thống.
\end{itemize}

\subsection{Kiểm thử}
Ghi nhận kết quả thực tế cho các kịch bản kiểm thử:
\begin{itemize}
  \item Sao chép thành công nội dung tệp tin nguồn sang tệp tin đích.
  \item Kiểm tra thông báo lỗi khi thiếu đối số dòng lệnh.
  \item Kiểm tra thông báo lỗi khi tệp nguồn không tồn tại.
\end{itemize}

\noindent Kết quả kiểm thử thực tế trên shell xv6:

\begin{center}
\begin{minipage}{0.86\textwidth}
\begin{lstlisting}[numbers=none,xleftmargin=0pt,framexleftmargin=0pt]
$ echo hello > a
$ cp a b
$ cat b
hello
$ cp a
usage: cp src dst
$ cp nofile out
cp: cannot open nofile
$
\end{lstlisting}
\end{minipage}
\end{center}
